% Options for packages loaded elsewhere
\PassOptionsToPackage{unicode}{hyperref}
\PassOptionsToPackage{hyphens}{url}
\documentclass[
]{article}
\usepackage{xcolor}
\usepackage{amsmath,amssymb}
\setcounter{secnumdepth}{-\maxdimen} % remove section numbering
\usepackage{iftex}
\ifPDFTeX
  \usepackage[T1]{fontenc}
  \usepackage[utf8]{inputenc}
  \usepackage{textcomp} % provide euro and other symbols
\else % if luatex or xetex
  \usepackage{unicode-math} % this also loads fontspec
  \defaultfontfeatures{Scale=MatchLowercase}
  \defaultfontfeatures[\rmfamily]{Ligatures=TeX,Scale=1}
\fi
\usepackage{lmodern}
\ifPDFTeX\else
  % xetex/luatex font selection
\fi
% Use upquote if available, for straight quotes in verbatim environments
\IfFileExists{upquote.sty}{\usepackage{upquote}}{}
\IfFileExists{microtype.sty}{% use microtype if available
  \usepackage[]{microtype}
  \UseMicrotypeSet[protrusion]{basicmath} % disable protrusion for tt fonts
}{}
\makeatletter
\@ifundefined{KOMAClassName}{% if non-KOMA class
  \IfFileExists{parskip.sty}{%
    \usepackage{parskip}
  }{% else
    \setlength{\parindent}{0pt}
    \setlength{\parskip}{6pt plus 2pt minus 1pt}}
}{% if KOMA class
  \KOMAoptions{parskip=half}}
\makeatother
\ifLuaTeX
  \usepackage{luacolor}
  \usepackage[soul]{lua-ul}
\else
  \usepackage{soul}
\fi
\setlength{\emergencystretch}{3em} % prevent overfull lines
\providecommand{\tightlist}{%
  \setlength{\itemsep}{0pt}\setlength{\parskip}{0pt}}
\ifPDFTeX
  \TeXXeTstate=1
  \newcommand{\RL}[1]{\beginR #1\endR}
  \newcommand{\LR}[1]{\beginL #1\endL}
  \newenvironment{RTL}{\beginR}{\endR}
  \newenvironment{LTR}{\beginL}{\endL}
\fi
\usepackage{bookmark}
\IfFileExists{xurl.sty}{\usepackage{xurl}}{} % add URL line breaks if available
\urlstyle{same}
\hypersetup{
  pdftitle={Jonah: },
  hidelinks,
  pdfcreator={LaTeX via pandoc}}

\title{Jonah:}
\usepackage{etoolbox}
\makeatletter
\providecommand{\subtitle}[1]{% add subtitle to \maketitle
  \apptocmd{\@title}{\par {\large #1 \par}}{}{}
}
\makeatother
\subtitle{(Whose Story Nevertheless Points to Christ)}
\author{}
\date{}

\begin{document}
\maketitle

Explanation of Satire

Satire uses folly or ridicule in order to expose wrongs and correct
them. As satire always has a sting, we should examine the story to see
where we find ourselves in the story. In this book we see Jonah as the
reluctant prophet who also stands for God's people as a whole. Look for
many inversions, subversions, and aversions as we go through the story.
Jonah also uses many examples of wordplay to show the ridiculousness of
the situation, such as fleeing from God, God's people being the only
creatures to not obey their creator, etc.

Explanation of this Guide

The text of Jonah is written in Optima while my notes are in Hoefler
Text. The paragraphs between the Biblical text give some commentary
while the footnotes go more in depth, including cross-references.

Outline of Jonah

1 Jonah's Descent\\
1.1 Jonah's Call and Rejection\\
1.2 Jonah at Sea in the Storm\\
1.3 Jonah Prays from the Deep\\
2 Jonah in a Foreign Land.\\
2.1 Jonah in Nineveh\\
2.2 Jonah Under the Gourd

Sources

I once told my daughter that often the bibliography is one of the most
interesting parts of a book. She responded `Then you must read boring
books'. To which I clarified that if I liked a work or wanted to get
more information regarding an author's argument the bibliography would
be the place to find more information. Some sources I used are:

The text of Jonah is mostly from the Lexham English Bible (LEB) and
Lexham English Septuagint (LES), altered with insight from the English
Standard Bible (ESV) and the New English Translation of the Septuagint
(NETS). The notes of the New English Translation (NET) (not related to
the NETS) were helpful in my footnotes.

Ancient Christian Commentary on Scripture: Old Testament XIV: The Twelve
Prophets

St Augustine, \emph{City of God}

Jonathan Pageau has an excellent article
(\href{https://orthodoxartsjournal.org/jonah/}{\ul{https://orthodoxartsjournal.org/jonah/}})
and YouTube video
(\href{https://www.youtube.com/watch?v=loL7O49iPDE}{\ul{https://www.youtube.com/watch?v=loL7O49iPDE}}
) on Jonah.

An overview on Jonah from the Bible Project:\\
\href{https://www.youtube.com/watch?v=dLIabZc0O4c}{\ul{https://www.youtube.com/watch?v=dLIabZc0O4c}}

Dictionary of Deities and Demons in the Bible (DDD)

1. Jonah's Decent

What to look for:\\
Jonah's descent, deeper and deeper, in these two chapters.\\
Where does he start and where does he end up? What are the steps along
the way?

\subsection{1.1 The Calling of Jonah, and Jonah's Rejection of that
Call}\label{the-calling-of-jonah-and-jonahs-rejection-of-that-call}

What to look for:\\
The similarities and differences between Jonah's call and the other
prophets.\\
Jonah's rejection of the call.\\
What is absent? Why did Jonah flee?

\textbf{Chapter 1~}And the Word of the LORD came\footnote{We should
  interpret this as the second person of the Trinity, not as some
  disembodied voice. See the call for other prophets but especially
  Samuel (1 Kingdoms/1 Samuel 3) where the Word of the Lord stands in
  front of the young Prophet Samuel.

  See
  \href{https://www.ancientfaith.com/podcasts/lordofspirits/the_word_of_the_lord/}{\ul{https://www.ancientfaith.com/podcasts/lordofspirits/the\_word\_of\_the\_lord/}}
  for more details.} to Jonah the son of Amittai\footnote{Jonah son of
  Amittai is referenced in 2 Kings (4 Kingdoms) 14.25 as a prophet to
  Northern Kingdom king Jeroboam II.}, saying, \textbf{2~}\RL{``}Get
up!\footnote{Literally `arise' which in addition to the contrast with
  Jonah getting up to go further and further downward is also is the
  word used for the `Resurrection'. Psalm 67/68.1 `Let God arise, and
  let his enemies be scattered' which we sing at Pascha.} Go to the
great city Nineveh and cry out\footnote{This phrase `cry out' is often
  used in the Old Testament (OT) in the context of injustice seeking to
  be remedied.} against her, because their evil has come up to my face
{[}gloss: come to my attention{]}.'' \textbf{3~}But Jonah got up to flee
toward Tarshish away from the face of the LORD. And he went down
\emph{to} Joppa and found a merchant ship going to Tarshish, and paid
her fare, and went on board her to go with them toward Tarshish away
from the face of the LORD.

This starts out in a typical fashion of the prophets who receive a
call.\footnote{For example, Isaiah 6, Jeremiah 1.4 ff, Ezekiel 2.3 ff.
  All variants of the calling of Moses in Exodus 3.} However, while most
prophets obey immediately or after a mild quibble, Jonah flat out
disobeys God and runs away from his presence toward the West\footnote{The
  West is the direction the sun sets and is a symbol of darkness. In our
  baptism we faced to the West and renounced the devil before turning
  toward the East and declaring we are uniting ourselves to Christ.
  Additionally, Christians pray facing the East and churches are built
  with the altars on the east end.}. And where is God's presence at this
time in history? Some of the Church Fathers say Jonah was in the Temple,
which might also be indicated by his prayer in chapter 2. So the
evil/disaster\footnote{The Hebrew word means both and the author
  playfully uses both meanings throughout the text.} of Nineveh has come
before the face of the Lord, the Lord tells Jonah, who is before his
faith, to get up cry out against them. (This call is unusual in itself
because the Lord usually sent his prophets to his people, but we see
that he has compassion on all people. His people were meant to be a
light to the nations.) Jonah initially gets up, appearing to obey this
call, but only to disobey and run from the face of God, as far as he can
to Tarshish in (what is now) Spain.

\subsection{1.2 Jonah at Sea in the
Storm}\label{jonah-at-sea-in-the-storm}

What to look for:\\
The repeated use of several words such as ``hurl'' and ``fear'' in the
description of the storm.

\textbf{4~}And the LORD hurled\footnote{In Greek `raised/aroused' but it
  is a different word than `arise'.} a great wind upon the sea, and it
was a great storm surge on the sea, and the merchant ship threatened to
break up.\footnote{In Hebrew the ship is personified, in that it
  considered shattering or threatened to shatter.} \textbf{5~}And the
mariners feared, and each cried out to his god\footnote{Wisdom 14.1-4
  ``Again {[}this is in a section about the uselessness of idols{]},
  someone preparing a voyage and about to travel through raging waves
  calls on wood more fragile than the ship bearing him. For desire for
  gain planned the ship but wisdom the craftsman made it. But your
  provision, Father, pilots the ship because you have made a way in the
  sea and as safe path in the waves.''}. And they hurled the
stuff\footnote{Usually translated `cargo' the word is more general like
  `objects'; its use here indicates the mariners were not picky about
  what they threw overboard, but whatever they could find. Of course
  there were also throwing over the cargo, indicating the storm was so
  severe they'd trade their livelihoods for their lives.} that was in
the merchant ship into the sea to lighten it for them. Meanwhile Jonah
went down into the belly\footnote{Usually `hold' but it's the same word
  used (in Greek) for the belly of the sea beast in 2.1 (1.17)} of the
vessel and lay down and fell asleep and was snoring.\footnote{The Greek
  states the snoring while the Hebrew states it was a deep sleep.}
\textbf{6~}And the captain of the ship approached him and said to him,
\RL{``}Why are you in a deep sleep?\footnote{``My soul, O my soul, rise
  up! Why art thou sleeping? The end draws near and soon thou shalt be
  troubled. Watch, then, that Christ thy God may spare thee, for He is
  everywhere present and fills all things.'' ---Kontakion of the Great
  Canon of Repentance.} Get up! Call on your god! Perhaps your god will
take notice of us and we won\RL{'}t perish!''

The sea is regarded as chaotic in the Bible (or even a manifestation of
chaos), it was unpredictable and full of monsters. But in the first
verse of this section we see God is in control. Like in the Genesis
story and Psalm 103, the Lord commands the waters and they obey. (There
is no battle between God and Chaos as in the pagan myths.) The pagan
mariners are shown to be more pious than the prophet as they call out to
their false gods, but he is sound asleep, oblivious to his own
destruction. Jonah went down into the hold (or `belly' in Greek) of the
ship, further showing his descent in the first half of the book. The
pagan captain of the ship, echos God's call to Jonah in the beginning of
the book, `Get up!' The comment `Perhaps your god will take
notice\ldots,' is echoed in 3.9 and will be discussed there.

During Lent, in the Great Canon of Repentance we call our souls to
awaken, which echoes the line here and in St Luke's Gospel (22.46) where
at the Mount of Olives, after Jesus prays before his crucifixion he
rebukes his disciples for sleeping.

This episode reminds us of the episode in St Mark's Gospel
(4.35--41)\footnote{Also Matthew 8.23--27 \& Luke 8.22--25 but the
  wordplay in Mark more closely echos Jonah.} where Jesus is asleep in
the stern of the boat in a great windstorm. His disciples call on him to
awake and do something. But of course Christ acts differently than
Jonah, as we discus below.

\textbf{7~}And the mariners said each to another, \RL{``}Come, let us
cast lots\footnote{``The lot is cast into the lap, but the outcome is of
  the Lord.'' Proverbs 16.33 (Hebrew). In the OT, lots are often cast to
  pick out a person. See Joshua 7.10--26, 1 Samuel (1 Kingdoms)
  10.20--24, 14.36--46, 2 Ezra 2.63, Luke 1.9, Acts 2.12--26, etc.} so
that we may know on whose account this disaster has come on us!'' And
they cast lots, and the lot fell on Jonah. \textbf{8~}So they said to
him, \RL{``}Please tell us whoever is responsible that this disaster has
come upon us! What is your occupation? And from where do you come? What
is your country? And from which people are you?'' \textbf{9~}And he said
to them, \RL{``}I am a Hebrew, a slave of the LORD\footnote{Hebrew: `I
  am a Hebrew', Greek: `I am a slave/servant of the Lord'.}, and I fear
the LORD, the God of heaven, who made the sea and the dry land.''
\textbf{10~}Then the men feared with great fear, and they said to him,
\RL{``}What is this you have done?'' because they knew that he was
fleeing from the face of the LORD (because he had told them).
\textbf{11~}So they said to him, \RL{``}What shall we do to you so that
the sea may quiet down for us?'' because the sea was growing more and
more tempestuous. \textbf{12~}And he said to them, \RL{``}Pick me up and
hurl me into the sea so that the sea may quiet down for you, because I
know that I am responsible that this great storm has come upon you
all.'' \textbf{13~}But the men rowed hard\footnote{Literally `digging
  {[}their oars{]} in'} to bring the ship back to the dry land, and they
could not do so because the sea was growing more and more tempestuous
against them. \textbf{14~}So they cried out to the LORD, and they said,
\RL{``}O LORD! Please do not let us perish because of this man\RL{'}s
life, and do not make us guilty of innocent blood, because you, O LORD,
did what you wanted.'' \textbf{15~}And they picked Jonah up and hurled
him into the sea, and the sea ceased from its raging. \textbf{16~}So the
men feared the LORD greatly, and they offered a sacrifice to the LORD
and made vows.

The pagan mariners recognized the Lord's hand in the lots and inquire of
Jonah. Jonah, who doesn't want to preach against the evil/disaster that
will befall Nineveh, brings evil/disaster upon the ship. He answers he's
a slave to the God of heaven who made both the sea and the dry land.
While this is a common truth expressed in the Scriptures it is comedic
in the mouth of Jonah, for it marks him as a runaway slave. But where is
he running to? He confesses God made everything, the \textbf{sea} and
the \textbf{dry land!} Psalm 138/139 states ``Where should I go from
your Spirit, and where would I flee from your face? If I should go up to
the Heavens, you are there. If I should go down to Hades/Sheol, you are
present. If I were to take up my wings at dawn, and dwell at the ends of
the sea, indeed, even there your hand will guide me and your right hand
will hold me.''

Does Jonah truly fear the LORD as he says? If the `fear of the LORD is
the beginning of wisdom' (Proverbs 9.10) is Jonah living wisely?

The sailors don't want to be guilty of Jonah's life and are now
beginning to fear the Lord (v. 10) so they row hard to get back to dry
land, but we are again told that the sea was growing more and more
tempestuous. So they listen to Jonah and hurl him overboard into the
deep. And immediately the sea was calmed and the mariners feared and
worshiped the true God.

In the Markan passage (4.35--41) Christ commands the storm to be still,
and it is instantly calmed. The disciples imitate the mariners here and
are filled with great fear, wondering who it is that can command the
winds and waves.

\subsection{1.3 Jonah Prays from the
Deep}\label{jonah-prays-from-the-deep}

What to look for:\\
The culmination of Jonah's descent and his first steps at repentance.

\textbf{Chapter 2} (\textbf{1.17)}\footnote{In most English Bibles this
  verse is numbered 1.17, but in Hebrew and Greek Bibles it's 2.1.}\textbf{~}And
the LORD appointed a large sea monster\footnote{Hebrew, ``great fish'',
  Greek ``great sea monster'' (sometimes rendered ``whale'' as in later
  Greek; the term came to mean ``whale''. But this is a representation
  of chaos.} to swallow up Jonah, and Jonah was in the belly of the
beast three days and three nights.\footnote{Hosea 6.1-2 ``Come, let us
  return to the LORD; because it is he who has torn, and he will heal
  us; he has struck us down and will bind us up. He will revive us after
  two days; on the third day he will raise us up, that we may live in
  his presence {[}literally: before his face{]}.}\textbf{~}

Just as the sea was considered chaotic, creatures in the deep were
manifestations of this chaos. Leviathan (the sea dragon) was chief among
them. However, here we see that unlike in pagan mythologies, God does
not struggle with chaotic creatures but commands or appoints them. Even
the chaos is subject to him. Psalm 103/104.26 states God formed the sea
dragon/leviathan and is played with by him.\footnote{Job 40-41
  (depending on the translation) also describes a fearsome leviathan/sea
  serpent. While often people try to identify leviathan with a
  real-world animal like a crocodile, plesiosaur, or colossal squid,
  this is a modernist mindset. It is more keeping with the text and
  ancient mindset to understand leviathan as an archetype of chaos and
  not associate it with a mundane animal.}

Jonah's three-day sojourn in the sea serpent is taken up by Christ, when
he advises an evil and adulterous generation that the only sign they
will be given is that of Jonah (Matthew 12.38-40, Luke 11.29-30). Just
as Jonah was three days in the deep, in the belly of the monster, in the
belly of Hades, so Christ was three days in the belly of the earth in
the tomb. But Christ harrowed (distressed) Hades and released all the
righteous from its clutches.

The prayer that follows is the basis for Ode 6 in the Matins/Orthros
canon.

\textbf{2} And Jonah prayed to the LORD his God from the belly of the
sea monster\footnote{Psalm 129/130 A song of ascents.

  Out of the depths I call to you, O LORD. \RL{``}O Lord, hear my voice.
  Let your ears be attentive to the voice of my supplications. If you, O
  LORD should keep track of iniquities, O Lord, who could stand?

  But with you is forgiveness, so that you may be feared. I await the
  LORD; my soul awaits, and I wait for his word. My soul waits for the
  Lord more than watchmen for the morning. Yes, more than watchmen for
  the morning.

  O Israel, hope for the LORD. For with LORD there is loyal
  {[}steadfast{]} love, and with him there is abundant redemption. And
  he will redeem Israel from all its iniquities.''} \textbf{3~}and said,

\RL{``}I called from my distress to the LORD,\\
\strut ~~~~and he answered me;\\
from the belly of Sheol\footnote{Hebrew, ``Sheol'', and Greek ``Hades''
  also sometimes translated ``Grave'' are the same place/state, the
  realm of the dead, conceived as being under the earth.} I cried for
help---\\
\strut ~~~~you heard my voice.\\
\textbf{4} And you threw me into the deep,\\
\strut ~~~~into the heart of the seas,\\
\strut ~~~~and the sea currents surrounded me;\\
all your breakers and your surging waves\footnote{Psalm 41.8/42.7 ``Deep
  calls to deep at the sound of your waterfalls. All your breakers and
  waves have passed over me.''}\\
\strut ~~~~passed over me.\\
\textbf{5} And I said, \RL{`}I am banished\\
\strut ~~~~from your sight;\\
how will I continue to look\\
\strut ~~~~on your holy temple?\RL{'}\\
\textbf{6}~The waters encompassed me up to my neck;\\
\strut ~~~~the deep surrounded me;\\
\strut ~~~~seaweed was wrapped around my head.\\
\textbf{7}~I went down to the foundations of the mountains;\\
\strut ~~~~the Underworld---its bars were around me
forever.\footnote{Psalm 17.5-7/18.4-6 ``The ropes/pains of death
  encircled me, and torrents of ruin overwhelmed me. The ropes of Sheol
  surrounded me; the snares of death confronted me.~In my trouble I
  called on the LORD and to my God I cried for help. He heard my voice
  from his temple, and my cry for help reached his ears.''}\\
But you brought up my life from the pit,\\
\strut ~~~~O LORD my God.\\
\textbf{8}~When my life was ebbing away from me,\\
\strut ~~~~I remembered the LORD,\\
and my prayer came to you,\\
\strut ~~~~to your holy temple.\\
\textbf{9}~Those who worship vain idols\\
\strut ~~~~forsake their loyal love.\\
\textbf{10}~But I, with a voice of thanksgiving,\\
\strut ~~~~will sacrifice to you;\\
I will fulfill what I have vowed.\\
\strut ~~~~Salvation belongs to the LORD!''

In verse 3 we see that Jonah has continued his descent, from the belly
of the ship to the belly of the fish to the belly of Hades. And in this
depth he prays to the Lord. And because the Lord is everywhere present
(Psalm 138), he hears Jonah. Notice this prayer of Jonah resonates with
many of the Psalms, some of which are noted in the footnotes. In verse 5
Jonah realizes how far he has fallen, from the presence of the Lord in
the temple to the depths of Hades. He continues this theme through verse
8, poetically stating he was bound fast in seaweed and the bars of the
Underworld held him in. In verse 7 he recognizes the Lord brought him up
from the pit, clearly he hadn't done anything to deserve it.

Verse 9 appears to be an indictment against the pagan mariners, but more
pointedly it's directed as those in Israel who forsake their God to
worship the gods of the nations. The pagans had not experienced the
steadfast love of the true God.

The last line of the prayer indicates that Jonah recognizes
salvation/deliverance comes from God, but we'll see how he reacts when
this is given to the Ninevites.

As a practical application, note that no matter how far we have fallen
from God, he is near to us. We can always call upon him, for as the
hymns state, he is the lover of mankind.

\textbf{11~}And the LORD spoke to the fish, and it vomited Jonah out on
the dry land.

Much as the Lord appointed the fish initially, he merely has to speak
and it obeys. In this first half we see God's creatures: the wind, the
waves, even the pagan mariners, and now this beast obey God, but his
prophet does not.

2 Jonah in a Foreign Land

What to look for:\\
Does Jonah ascend as he once descended?\\
What is his attitude to the inhabitants of Nineveh?

\subsection{2.1 Jonah in Nineveh}\label{jonah-in-nineveh}

What to look for:\\
Parallels between chapters 3 \& 1\\
What is Jonah's message?\\
How is Jonah's message received?

\textbf{Chapter 3~}And the Word of the LORD came to Jonah a second time,
saying, \textbf{2~}\RL{``}Get up! Go to Nineveh, the great city, and
proclaim to it the message that I am telling you.'' \textbf{3~}So Jonah
got up and went to Nineveh according to the Word of the LORD.

Here we see the Lord is gracious and again speaks his word to Jonah.
This time Jonah gets up and goes to Nineveh. The text sets up an
expectation of Jonah's ascent. His ascent in this second part would
correspond to his descent in the first two chapters. The watery
narrative also calls to mind the Flood of Noah's day which unmade the
world, but the world was remade after the waters receded. Will we find a
similar arc with Jonah?

Now Nineveh was an extraordinarily great city---a journey of three
days.\footnote{There are various ways interpreters explain this phrase.
  Let's say that it idiomatically indicates Nineveh was an important
  city. Note the parallel with other parts of the book mentioning three.}
\textbf{4~}And Jonah began to go into the city a journey of one day, and
he cried out and said, \RL{``}Forty\footnote{Greek `three', Hebrew
  `forty'.

  St Augustine explains how we may derive benefit from this seeming
  contradiction in The City of God: Book 18: Chapter 44:

  How the Threat of the Destruction of the Ninevites is to Be Understood
  Which in the Hebrew Extends to Forty Days, While in the Septuagint It
  is Contracted to Three.

  But some one may say, "How shall I know whether the prophet Jonah said
  to the Ninevites, Yet three days and Nineveh shall be
  overthrown,\textquotesingle{} or forty days?" {[}1219{]} For who does
  not see that the prophet could not say both, when he was sent to
  terrify the city by the threat of imminent ruin? For if its
  destruction was to take place on the third day, it certainly could not
  be on the fortieth; but if on the fortieth, then certainly not on the
  third. If, then, I am asked which of these Jonah may have said, I
  rather think what is read in the Hebrew, "Yet forty days and Nineveh
  shall be overthrown." Yet the Seventy, interpreting long afterward,
  could say what was different and yet pertinent to the matter, and
  agree in the self-same meaning, although under a different
  signification. And this may admonish the reader not to despise the
  authority of either, but to raise himself above the history, and
  search for those things which the history itself was written to set
  forth. These things, indeed, took place in the city of Nineveh, but
  they also signified something else too great to apply to that city;
  just as, when it happened that the prophet himself was three days in
  the whale\textquotesingle s belly, it signified besides, that He who
  is Lord of all the prophets should be three days in the depths of
  hell. Wherefore, if that city is rightly held as prophetically
  representing the Church of the Gentiles, to wit, as brought down by
  penitence, so as no longer to be what it had been, since this was done
  by Christ in the Church of the Gentiles, which Nineveh represented,
  Christ Himself was signified both by the forty and by the three days:
  by the forty, because He spent that number of days with His disciples
  after the resurrection, and then ascended into heaven, but by the
  three days, because He rose on the third day. So that, if the reader
  desires nothing else than to adhere to the history of events, he may
  be aroused from his sleep by the Septuagint interpreters, as well as
  the prophets, to search into the depth of the prophecy, as if they had
  said, In the forty days seek Him in whom thou mayest also find the
  three days,-\/-the one thou wilt find in His ascension, the other in
  His resurrection. Because that which could be most suitably signified
  by both numbers, of which one is used by Jonah the prophet, the other
  by the prophecy of the Septuagint version, the one and self-same
  Spirit hath spoken. I dread prolixity, so that I must not demonstrate
  this by many instances in which the seventy interpreters may be
  thought to differ from the Hebrew, and yet, when well understood, are
  found to agree. For which reason I also, according to my capacity,
  following the footsteps of the apostles, who themselves have quoted
  prophetic testimonies from both, that is, from the Hebrew and the
  Septuagint, have thought that both should be used as authoritative,
  since both are one, and divine.} more days and Nineveh will be
overthrown!'' \textbf{5~}And the people of Nineveh believed in God, and
they proclaimed a fast and put on sackcloth---from the greatest of them
to the least.

\textbf{6~}And the news reached the king of Nineveh, and he rose from
his throne and removed his royal robe, put on sackcloth, and sat in the
ashes.

Many Church fathers note that the message of God came to the
nations/gentiles only after Jonah was in the belly of the sea monster 3
days and nights, just like after Christ was in the earth the message of
the Gospel was preached to the nations.

According to the text, Jonah\RL{'}s message was only \RL{`}Forty days
and Nineveh will be overthrown.\RL{'} No message of repentance or that
they could avoid this judgement. But even with this little information,
like the pagan mariners in Jonah 1, the pagan Ninevites believe in God
and repent. In the Gospels, Jesus says the generation of Nineveh will
rise up and judge those who had more Light and did not follow it
(Matthew 12.41 \& Luke 13.32). Given that Nineveh repents, there is some
humorous irony here that Jonah is made to look like a false prophet.

The Oxford English Study Bible notes an interesting play on words with
`overthrown' stating that in addition to the obvious meaning it can also
mean `turned about' or `turned around'. As such the repentance of the
city involved an upheaval of their wicked way of life and their lives
being turned around to live before the face of God.

\textbf{7~}And he had a proclamation made, and said, \RL{``}In Nineveh,
by a decree of the king and his nobles: \RL{``}No human being or animal,
no herd or flock, shall taste anything! They must not eat, and they must
not drink water! \textbf{8~}And the human beings and the animals must be
covered with sackcloth! And they must call forcefully to God, and each
must turn from his evil way and from the violence that is in his hands.
\textbf{9~}Who knows? God may relent and change his mind and turn from
his blazing anger so that we will not perish.'' \textbf{10~}And God saw
their deeds---that they turned from their evil ways---and God changed
his mind about the evil that he had said he would bring upon them, and
he did not do it.

In the satirical, exaggerated style of Jonah, even the cattle are called
on to fast and wear sackcloth as part of the repentance of Nineveh. The
last sentence echos what the mariners said as they were facing
destruction in the deep and reflects Joel 2.14\footnote{As far back as
  we have records all the Minor Prophets have been bound together on one
  scroll/book---The Twelve. They have been received as a unit and should
  be studied as one, but we don't have the space to explore this in
  depth here.}. In the book of Joel, God has sent an army of locusts to
devastate the land in judgment. At the end of that plague, the Lord
calls his people to repentance ``rend your hearts and not your
garments''\footnote{That is, perform true repentance and not merely go
  through the outward signs.} and the prophet states, ``Who knows
whether he will turn and relent, and leave a blessing behind him, an
offering and a libation, for the LORD your God?'' This is reflected in
our Liturgy where the priest lifts up the offering saying, ``Thine own
of thine own, we offer unto thee in behalf of all and for all.'' Even
our ability to worship God comes from him.

In Jeremiah 18, we see that if the Lord declares judgment on a nation
but they repent and turn from their evil he will relent of the
calamity\footnote{Remembering `evil' and `calamity' are the same word in
  Hebrew.} he intended. St John Chrysostom uses this passage to point
out even as evil as Nineveh was, God is exceedingly gracious and
accepted their repentance after only three days.

\subsection{2.2 Jonah Under the Gourd}\label{jonah-under-the-gourd}

What to look for:\\
Parallels (anti-parallels) between chapters 4 \& 2\\
Contrast between God and Jonah.

\textbf{Chapter 4~}And this was evil to Jonah, a great evil\footnote{This
  is the literal translation from Hebrew, which continues the
  evil/calamity wordplay throughout the book. Or `exceedingly evil', or
  `displeasing with great displeasure' or from the Geek, `Jonah was
  grieved with great grief''.}, and he became blazingly
angry.\footnote{Literally `it was hot for him'. Contrast with 3.9
  referring to God's anger, God's question to Jonah in 4.4 and the heat
  of the desert in 4.8-9. Greek `confounded' or `confused'.}
\textbf{2~}And he prayed to the LORD and said, \RL{``}O LORD, were these
not my words\footnote{Contrast with the Word of the Lord in 1.1 \& 3.1.}
while I was in my homeland? Therefore I originally fled\footnote{The
  Hebrew is difficult here. Alternately as in LXX `had the foresight to
  flee' or `anticipated to flee'.} to Tarshish\footnote{``Jonah fled to
  Tarshish, foreseeing the conversion of the men of Nineveh; for as a
  prophet he knew the loving-kindness of God, but he was jealous that
  his prophecy should not be proved false\ldots.All the names of the Old
  Testament have I set before thee, my soul, as an example. Imitate the
  holy acts of the righteous and flee from the sins of the wicked.''
  ---From Ode 8 of the Canon of Repentance.}, because I knew that you
are a gracious and compassionate God, slow to anger and having great
steadfast love, and one who relents concerning calamity\footnote{Joel
  2.13: Rend your hearts and not your garments, and return to the LORD
  your God, because he is gracious and compassionate, slow to anger and
  great in loyal love, and relenting from harm.}. \textbf{3~}And so
then, O LORD, please take my life from me, because for me death is
better than life!'' \textbf{4~}And the LORD said, \RL{``}Is it right for
you to be blazingly angry?''

Now we hear why Jonah ran from the Lord, he hated---it was exceedingly
evil to him---that the Lord was compassionate and forgiving. While Jonah
was fine with the Lord saving him from the depths of hades in chapter 2,
when he does to the great city of Nineveh he is grieved. He also quotes
Joel, (2.13) the verse immediately preceding what the King of Nineveh
quoted, but he throws it in God's face. While God is concerned about the
nations, Jonah only wants salvation for his people.\footnote{There is a
  contrast here between Jonah who flees Nineveh when God spares it and
  Elijah who flees to the wilderness when the Queen Jezebel threatens
  his life (1 Kings {[}3 Kingdoms{]} 19.4--5).}

\textbf{5~}And Jonah went out from the city and sat down east of the
city,\footnote{Genesis 4.16 ``And Cain went out from the presence of the
  LORD, and he settled in the land of Wandering, east of Eden.'' And
  Genesis 2.8 ``And the LORD God planted a garden in Eden to the east,
  and there he put the man he had formed.''} and he made for himself a
shelter there. And he sat under it in the shade, waiting to see what
would happen with the city. \textbf{6~}And the LORD God appointed a
plant, and he made it grow up over Jonah to be a shade over his head, to
save\footnote{``shade'' and ``save'' sound similar in Hebrew. Also
  `misery' can be rendered `discomfort', or `evil' as the NET Bible
  notes indicate, there is a lot of `polysemantic wordplay' going on
  here.} him from his misery. And Jonah rejoiced with great
joy\footnote{Note the contrast to 4.1.} about the plant. \textbf{7~}So
God appointed a worm at daybreak the next day, and it attacked the
plant, and it withered. \textbf{8~}And when the sun rose, God appointed
a scorching\footnote{The scorching wind paralleling Jonah's anger. But
  also possibly a reference to Genesis 3.8 with the Lord coming in the
  Spirit/wind of the Day (usually rendered `cool of the day' for obscure
  reasons we don't have the space to address here).} east wind, and the
sun attacked Jonah\RL{'}s head and he grew faint. And he asked that he
could die and said, \RL{``}My death is better than my life!''
\textbf{9~}So God said to Jonah, \RL{``}Is it right for you to be
blazingly angry about the plant?'' And he said, \RL{``}It is right for
me to be blazingly angry enough to die!''

While Jonah sinks into his hot anger, the Lord seeks to save him from
his misery. Parallel to him appointing a sea monster in chapter 2, God
appoints a plant to grow overnight to shade Jonah. But as God appoints a
worm to smite the plant and appoints the hot east wind and the sun to
smite Jonah. We see the pattern continuing of everything in God's
creation obeying God without complaint except his prophet. Let us who
are the people of God take note and not emulate Jonah.

\textbf{10~}But the LORD said, \RL{``}You are troubled about the plant,
for which you did not labor nor cause it to grow. It grew up in a night
and it perished in a night! \textbf{11~}And should I not be concerned
about Nineveh, the great city, in which there are more than one hundred
and twenty thousand people who do not know right from left, plus many
animals?''

How gracious is the Lord, in that he confronts Jonah in his sin! The
Lord tells Jonah it is bizarre for him to be troubled about a plant that
he did not work for and yet not care at all about the great city of
Nineveh, in which 120,000 people don't know left from right, up from
down, or good from evil. The concluding statement (plus much cattle)
plays on Jonah's concern for the plant, perhaps he'd also be concerned
about the cattle that would perish with the city, after all they fasted
and put on sackcloth too.

So is Jonah a saint in the Orthodox Church? Why is he painted on our
dome?

\RL{``}Now what do you think? A man had two sons. He approached the
first and said, \RL{`}Son, go work in the vineyard today.\RL{'} And he
answered and said, \RL{`}I do not want to!\RL{'} But later he changed
his mind and went. And he approached the second and said the same thing.
So he answered and said, \RL{`}I will, sir,\RL{'} and he did not go.
Which of the two did the will of his father?'' They said, \RL{``}The
first.'' Jesus said to them, \RL{``}Truly I say to you that the tax
collectors and the prostitutes are going ahead of you into the kingdom
of God!~For John came to you in the way of righteousness and you did not
believe him, but the tax collectors and the prostitutes did believe him.
And when you saw it, you did not even change your minds later so as to
believe in him. ---Mathew 21.28--32

Jonah prefigures the death and resurrection of Christ, and traditionally
wrote the book with his name, showing his understanding of his situation
in hindsight.

But also where do we see ourselves in the Book of Jonah? Are we the
prophet, with correct knowledge of God but not his compassion. Are we
the Ninevites who quickly repent on hearing the Gospel?

\emph{Through the prayers of the Prophet Jonah, O Lord Jesus Christ, our
God, have mercy upon us and save us.}

\end{document}
